\section{问题定义}
\subsection{复杂地形下的人机牵引柔顺运动控制}

人通过机器人手腕、把手或机身施加牵引力$F^s$，引导机器人平稳移动与转向，并在持续牵引下自适应通过复杂地形、保持动态平衡；牵引解除后，机器人平滑恢复原始速度指令。

\subsection{网络输入与输出}
每帧观测为 45 维，按以下顺序拼接：
\begin{equation}
\bm o_t=\left[
\bm c_t^{\rm obs},\;
\bm\omega_t^{\rm b},\;
\bm g_t^{\rm b},\;
(\bm q_t-\bm q_{\rm nom}),\;
\dot{\bm q}_t,\;
\bm a_{t-1}
\right]^\top\in\R^{45}.
\label{eq:proprio-observation}
\end{equation}
六项依次为速度指令（3）、机身角速度（3）、重力投影（3）、相对关节位置（12）、关节速度（12）和上一次关节动作（12）。其中 $\bm c_t^{\rm obs}$ 为观测中的速度指令。最近 $H$ 帧组成网络输入
\begin{equation}
h_t=\left[\bm o_t,\;\bm o_{t-1},\;\ldots,\;\bm o_{t-H+1}\right]^\top\in\R^{45H}.
\label{eq:proprio-history}
\end{equation}
两套方案均使用 $H=6$。前向、侧向和偏航是速度接口的三个分量。
旧版 H6 的网络输出速度修正，纯牵引策略直接输出速度指令。
各自的网络输入和反馈时序见第~\ref{sec:h6-baseline} 和第~\ref{sec:pure-traction} 节。

\subsection{速度指令与牵引力转换}

速度指令的平移分量采用 $\mathrm{m/s}$，偏航分量采用 $\mathrm{rad/s}$。

牵引力作用于点 $P$ 时，将力转换到机身坐标系 $b$，并计算关于机身参考点的等效力矩。
\begin{equation}
\bm F^b=\bm R_{bs}\bm F^s,
\qquad \bm\tau_{\rm eq}^b=\bm r_{bP}^b\times\bm F^b.
\label{eq:force-transform}
\end{equation}
其中 $\bm R_{bs}$ 将 $s$ 系向量旋转到 $b$ 系，$\bm r_{bP}^b$ 是机身参考点指向 $P$ 的位矢，以 $b$ 系表示。
